% ============================================================================
% Fichier  : partie1/P01_C03_evaluer_controle_interne.tex
% Titre    : Évaluer le Système de Contrôle Interne
% Auteur   : MINKA MI NGUIDJOÏ Thierry Emmanuel
% Version  : 0.1 -- novembre 2026
% ============================================================================
% Chapitre central du cours (décision de l'auteur, 07/11/2026).
% Sources : COSO 2013 (IFACI), chap. 3, §§ 138-140 (déficience, déficience
% majeure) ; INF3226 C03 (tailles d'échantillon, NON / NEK) ; décret
% n° 2013/287, art. 28.
% Faits MINÉQUIR établis ici (à garder) :
%  - contrôle 5 du P00 : pièces du poste comptable reçues le 2 avril 2026,
%    le contrôle des pièces jointes a fonctionné -> OUI ; lecture finale
%    3 OUI, 5 NON, 0 NEK ;
%  - 31 mutations ou départs en 2025 ; 9 comptes encore actifs au 23/03/2026 ;
%  - déficiences : A (accès et séparation, contrôles 1 et 2) majeure ;
%    B (fichier des fournisseurs et doublons, contrôles 3 et 4) ;
%    C (revue de fin de gestion, contrôle 7).
% ============================================================================

\chapter{Évaluer le Système de Contrôle Interne}
\label{chap:eval-sci}
\chaptermark{Évaluer le contrôle interne}

\epigraph{%
    « Un contrôle mal conçu est inefficace, même exécuté à la lettre~; un
    contrôle bien conçu mais non exécuté ne protège de rien. »%
}{Règle du cours}

\niveauChapitre{Intermédiaire}{Chapitre~\ref{chap:sci}.}

\begin{boxobjectifs}
\begin{itemize}
    \item Conduire l'évaluation d'un système de contrôle interne en cinq
          étapes~: décrire, questionner, cheminer, évaluer la conception,
          tester le fonctionnement.
    \item Construire une matrice des risques et des contrôles pour un
          circuit financier informatisé.
    \item Utiliser un questionnaire de contrôle interne, en distinguant la
          réponse de l'audité de la preuve.
    \item Réaliser un test de cheminement et en tirer ce qu'il établit, et
          ce qu'il n'établit pas.
    \item Dimensionner un test de fonctionnement et coter le contrôle
          \OUI{}, \NON{} ou \NEK{}.
    \item Classer les déficiences (déficience, déficience majeure) selon le
          COSO~2013.
    \item Formuler une conclusion sur l'efficacité du système de contrôle
          interne.
\end{itemize}
\end{boxobjectifs}

\bigskip

% ============================================================================
\section*{Pourquoi ce chapitre}
% ============================================================================

Dans le prologue, la vérificatrice a coté huit contrôles. Elle savait
quels contrôles attendre, comment les tester, et ce que voulait dire un
\NON{}. Ce chapitre donne cette démarche, pas à pas. Il est le plus
important du cours~: chaque chapitre de la Partie~IV l'applique à un
domaine, et le cas intégrateur la reprend en entier.

La question directrice est celle de l'article~28 du décret
n°~2013/287~: évaluer les structures de contrôle interne. Elle se
décompose en deux questions plus précises. \emph{Le dispositif est-il
conçu pour réduire les risques~?} \emph{A-t-il fonctionné pendant toute
la période~?}

\separateur

% ============================================================================
\section{La démarche en cinq étapes}
\label{sec:c03-demarche}
% ============================================================================

\begin{tableaucours}{cL{2.6cm}L{4.8cm}L{4.4cm}}{%
    \tableauHeader{N°} &
    \tableauHeader{Étape} &
    \tableauHeader{Question} &
    \tableauHeader{Production}
}
1 & Décrire & Comment le circuit fonctionne-t-il, et où sont les
risques~? & Narratif, diagramme, matrice des risques et des contrôles \\
2 & Questionner & Quels contrôles l'entité dit-elle exercer~? &
Questionnaire de contrôle interne rempli \\
3 & Cheminer & Le circuit fonctionne-t-il comme on me l'a décrit~? &
Fiche de test de cheminement \\
4 & Évaluer la conception & Ces contrôles, s'ils fonctionnent,
réduisent-ils les risques~? & Cotation de la conception \\
5 & Tester le fonctionnement & Ont-ils fonctionné toute la période~? &
Feuilles de test, cotation finale \\
\end{tableaucours}
\vspace{-0.8em}
\begin{center}{\small\itshape Tableau~3.1. Les cinq étapes de l'évaluation
du contrôle interne}\end{center}

Les étapes s'enchaînent, et chacune peut arrêter la suivante. Un contrôle
dont la conception est défaillante n'a pas à être testé dans son
fonctionnement~: il est déjà une déficience. Un contrôle que l'entité ne
déclare pas, mais que la matrice des risques exige, est un \NON{} dès
l'étape~2. Les étapes~6 et~7, classer les déficiences et conclure, sont
traitées aux sections~\ref{sec:c03-deficiences} et~\ref{sec:c03-conclure}.

\separateur

% ============================================================================
\section{Décrire le système}
\label{sec:c03-decrire}
% ============================================================================

\subsection{Le narratif et le diagramme}

L'auditeur commence par décrire le circuit, avec les mots de l'entité~:
qui fait quoi, dans quel ordre, avec quel écran, et quelle trace reste.
Le narratif s'accompagne d'un diagramme, qui fait apparaître les étapes,
les acteurs et les points de contrôle.

\begin{center}
\begin{tikzpicture}[
    etape/.style = {draw=CouleurBordDef, fill=CouleurDef, rounded corners=3pt,
        text width=2.5cm, minimum height=0.9cm, align=center,
        font=\scriptsize\bfseries\color{CouleurPrimaire}},
    ctrl/.style = {draw=CouleurAccent, fill=CouleurAccent!15, circle,
        inner sep=1pt, font=\scriptsize\bfseries},
    arrow/.style = {-{Stealth[length=4pt]}, thick, color=CouleurSecondaire}
]
\node[etape] (e) at (0,0) {Engagement\\(ENG)};
\node[etape] (l) at (3.4,0) {Liquidation\\(LIQ)};
\node[etape] (o) at (6.8,0) {Ordonnancement\\(ORD)};
\node[etape] (p) at (10.2,0) {Paiement\\(poste comptable)};
\node[etape, text width=3.4cm] (f) at (5.1,-2.0) {Fichier des fournisseurs\\(FRS)};
\draw[arrow] (e) -- (l);
\draw[arrow] (l) -- (o);
\draw[arrow] (o) -- (p);
\draw[arrow, dashed] (f) -- (l);
\draw[arrow, dashed] (f) -- (o);
\node[ctrl] at (1.7,0.75) {C1};
\node[ctrl] at (5.1,0.75) {C2};
\node[ctrl] at (8.5,0.75) {C3};
\node[ctrl] at (7.4,-2.0) {C4};
\end{tikzpicture}
\end{center}
\vspace{-0.4em}
\begin{center}{\small\itshape Figure~3.1. Le circuit de la dépense dans
SIDÉPENSE et ses points de contrôle (C1 contrôle du crédit~; C2 contrôle
de la liquidation~; C3 contrôle des doublons et des pièces~; C4 contrôle
des modifications de RIB)}\end{center}

\subsection{La matrice des risques et des contrôles}

La description aboutit à une \termecle{matrice des risques et des
contrôles}\index{matrice des risques et des contrôles}~: pour chaque
risque du circuit, le ou les contrôles qui doivent le couvrir. C'est la
liste de ce que l'auditeur s'attend à trouver, établie \emph{avant}
d'interroger l'entité, à partir des textes, des référentiels et de son
expérience.

\begin{tableaucours}{L{4.0cm}L{4.6cm}L{1.6cm}L{1.8cm}}{%
    \tableauHeader{Risque} &
    \tableauHeader{Contrôle attendu} &
    \tableauHeader{Type} &
    \tableauHeader{Principe}
}
R1. Un agent parti agit encore & Désactivation du compte à la date de la
mutation & Préventif & 11 \\
R2. Un agent émet seul un mandat & Profils LIQ et ORD jamais réunis &
Préventif & 10, 11 \\
R3. Le bénéficiaire est changé & Seconde validation de tout RIB modifié &
Préventif & 10 \\
R4. Une facture est payée deux fois & Blocage des doublons dans
l'application & Préventif & 10 \\
R5. Le mandat n'est pas justifié & Vérification des pièces jointes par le
poste comptable & Préventif & 10 \\
R6. Une opération ne laisse pas de trace & Journalisation des opérations
sensibles & Détectif & 13 \\
R7. Une anomalie de fin de gestion passe & Revue des mandats émis hors
heures et des RIB modifiés & Détectif & 16 \\
R8. Les traces disparaissent & Conservation protégée des journaux &
Détectif & 13 \\
\end{tableaucours}
\vspace{-0.8em}
\begin{center}{\small\itshape Tableau~3.2. Matrice des risques et des
contrôles du circuit de la dépense (MINÉQUIR)}\end{center}

\begin{boxalert}
\textbf{Violation typique.} Construire la liste des contrôles à partir de
ce que l'entité déclare faire. On ne trouve alors que des contrôles qui
existent~; ceux qui manquent n'apparaissent jamais.

\smallskip
\textbf{Règle~:} la matrice part des risques, pas des contrôles déclarés.
Un risque sans contrôle en face est déjà une constatation.
\end{boxalert}

\separateur

% ============================================================================
\section{Questionner : le questionnaire de contrôle interne}
\label{sec:c03-qci}
% ============================================================================

\begin{boxdef}
\termecle{Questionnaire de contrôle interne}\index{questionnaire de contrôle
interne} (\sigle{QCI})~: liste de questions fermées, tirée de la matrice des
risques et des contrôles, qui demande à l'entité si chaque contrôle attendu
existe, qui l'exerce, à quelle fréquence, et quelle trace il laisse.
\end{boxdef}

Une question bien formée appelle une réponse vérifiable. «~Les droits
sont-ils bien gérés~?~» n'appelle qu'un «~oui~». «~Qui désactive le compte
d'un agent muté, dans quel délai, et sur quelle pièce~?~» appelle une
réponse que l'on peut tester.

\begin{tableaucours}{cL{5.4cm}L{3.0cm}L{3.2cm}}{%
    \tableauHeader{N°} &
    \tableauHeader{Question} &
    \tableauHeader{Réponse déclarée} &
    \tableauHeader{Pièce demandée}
}
1 & Qui informe la cellule informatique d'une mutation, et sous quel
délai~? & «~Le service du personnel, en principe~» & Procédure, exemples
de notifications \\
2 & Qui désactive le compte, et quelle trace reste~? & «~La cellule, dès
qu'elle est informée~» & Journal des désactivations 2025 \\
3 & Les profils LIQ et ORD peuvent-ils être réunis sur un compte~? &
«~Non~» & Matrice des profils \\
4 & Qui peut modifier un RIB, et avec quelle validation~? & «~Le service
de la liquidation~» & Liste des détenteurs du profil FRS, journal des
modifications \\
5 & L'application bloque-t-elle les doublons~? & «~Oui, le numéro de
facture est unique~» & Paramétrage, essai \\
6 & Une revue des mandats de fin de gestion est-elle faite~? & «~Le DAF
suit les gros montants~» & Comptes rendus de revue \\
\end{tableaucours}
\vspace{-0.8em}
\begin{center}{\small\itshape Tableau~3.3. Extrait du QCI sur le circuit
de la dépense, rempli par la DAF et la cellule informatique}\end{center}

\begin{boxinfo}
\textbf{Une réponse au QCI est une déclaration, pas une preuve.} Elle dit
ce que l'entité croit faire. Le cheminement et les tests disent ce qu'elle
fait. La réponse à la question~5 en est l'exemple~: le numéro de facture
est bien unique, mais le suffixe «~bis~» suffit à en faire un autre
numéro. Le contrôle déclaré existe~; il ne couvre pas le risque.
\end{boxinfo}

\separateur

% ============================================================================
\section{Cheminer : le test de cheminement}
\label{sec:c03-cheminement}
% ============================================================================

\begin{boxdef}
\termecle{Test de cheminement}\index{test de cheminement}
(\textit{walkthrough})~: suivi d'une opération réelle, du début à la fin
du circuit, en examinant à chaque étape l'écran utilisé, la trace produite
et le contrôle exercé. Il établit que l'auditeur a compris le circuit et
que les contrôles existent tels qu'ils sont décrits~; il n'établit pas
qu'ils ont fonctionné pendant toute la période.
\end{boxdef}

Mme~NGO~BIYONG choisit le mandat 2025-08817, le premier paiement du
décompte n°~3, et le suit dans SIDÉPENSE, assise à côté de l'agent qui
l'a traité.

\begin{tableaucours}{L{2.4cm}L{5.4cm}L{4.4cm}}{%
    \tableauHeader{Étape} &
    \tableauHeader{Ce que la vérificatrice observe} &
    \tableauHeader{Trace obtenue}
}
Engagement & Le crédit est contrôlé~: un essai à un montant supérieur est
refusé & Capture de l'écran de refus \\
Liquidation & L'agent joint la facture et le PV de réception, puis valide
& Historique du mandat \\
Ordonnancement & Un autre agent valide~; l'application ne vérifie pas
qu'il est différent du liquidateur & Historique (deux noms différents,
par choix et non par contrainte) \\
Transmission & Le mandat part au poste comptable avec ses pièces
scannées & Bordereau électronique \\
Paiement & Le poste comptable vérifie les pièces, puis paie & Visa
électronique, relevé \\
\end{tableaucours}
\vspace{-0.8em}
\begin{center}{\small\itshape Tableau~3.4. Test de cheminement du mandat
2025-08817}\end{center}

\begin{boxscenario}{Dossier MINÉQUIR. Ce que montre le cheminement}
\textbf{Constatation.} Sur le mandat 2025-08817, liquidation et
ordonnancement ont été faits par deux agents différents. L'application
n'empêche pas qu'un même agent fasse les deux.

\medskip
\textbf{Analyse} (confiance élevée). La séparation a été respectée sur ce
mandat, mais par habitude, pas par contrainte. Un cheminement réussi ne
prouve donc pas que le contrôle existe~: il montre ici que la séparation
dépend de la bonne volonté des agents. C'est un défaut de
\emph{conception}, que le cheminement suffit à établir.
\end{boxscenario}

\separateur

% ============================================================================
\section{Évaluer la conception, tester le fonctionnement}
\label{sec:c03-tester}
% ============================================================================

\subsection{Deux questions distinctes}

\begin{tableaucours}{L{3.0cm}L{4.8cm}L{4.4cm}}{%
    \tableauHeader{Évaluation} &
    \tableauHeader{Question} &
    \tableauHeader{Comment}
}
Conception & Le contrôle, s'il fonctionne, réduit-il le risque~? &
Lecture de la procédure, du paramétrage, cheminement \\
Fonctionnement & Le contrôle a-t-il fonctionné comme prévu pendant toute
la période~? & Test d'un échantillon d'exécutions et de leurs traces,
ou test du paramétrage \\
\end{tableaucours}
\vspace{-0.8em}
\begin{center}{\small\itshape Tableau~3.5. Conception et fonctionnement
d'un contrôle}\end{center}

\subsection{Combien d'exécutions tester~?}

Le fonctionnement d'un contrôle manuel se teste sur un échantillon dont la
taille dépend de la fréquence du contrôle. Le cours retient les ordres de
grandeur du Tableau~3.6~; ce sont des conventions de travail, que la
méthodologie de chaque institution peut fixer autrement.

\begin{tableaucours}{L{3.6cm}L{3.2cm}L{5.4cm}}{%
    \tableauHeader{Fréquence du contrôle} &
    \tableauHeader{Exécutions par an} &
    \tableauHeader{Échantillon indicatif}
}
Annuelle & 1 & 1 \\
Trimestrielle & 4 & 2 \\
Mensuelle & 12 & 2 à 5 \\
Hebdomadaire & 52 & 5 à 15 \\
Quotidienne & environ 250 & 20 à 40 \\
Plusieurs fois par jour & plus de 250 & 25 à 60 \\
Automatisée & & Un essai qui doit passer et un qui doit être bloqué, si
les contrôles généraux sont efficaces \\
\end{tableaucours}
\vspace{-0.8em}
\begin{center}{\small\itshape Tableau~3.6. Tailles d'échantillon
indicatives (convention du cours)}\end{center}

\subsection{Coter}

\begin{itemize}
    \item \textbf{\OUI{}}~: le contrôle est bien conçu et a fonctionné sur
          tout l'échantillon (ou l'essai automatisé a donné le résultat
          attendu).
    \item \textbf{\NON{}}~: la conception est défaillante, ou une
          exécution au moins a manqué, ou la trace que le contrôle devait
          produire n'existe pas.
    \item \textbf{\NEK{}}~: malgré des diligences normales, l'auditeur ne
          peut pas conclure~: pièce non remise, sources contradictoires
          qu'il ne peut pas départager. Il écrit les diligences faites et
          relance.
\end{itemize}

\begin{boxobjet}{C-ACC-01 · Comptes des agents mutés}
\textbf{Objectif.} Aucun agent muté ou parti ne peut agir dans
SIDÉPENSE.

\textbf{Risque couvert.} R1~: un compte resté actif est utilisé par son
ancien titulaire ou par un tiers.

\textbf{Critère.} Le compte est désactivé au plus tard à la date d'effet
de la mutation (procédure à écrire~; à défaut, principe~11 du COSO).

\textbf{Test de conception.} Existe-t-il une notification du service du
personnel à la cellule informatique~? Réponse~: non, aucune procédure
écrite.

\textbf{Test de fonctionnement.} Rapprocher la liste des 31 mutations et
départs de 2025 de la liste des comptes actifs au 23~mars~2026.

\textbf{Résultat.} 9 comptes sur 31 sont encore actifs, dont
\texttt{mbarga.j}.

\textbf{Cote.} \NON{} (conception et fonctionnement défaillants).
\end{boxobjet}

\begin{boxalert}
\textbf{Violation typique.} Coter \NEK{} un contrôle dont la procédure
exige un procès-verbal, alors que l'entité a confirmé qu'il n'en existe
aucun.

\smallskip
\textbf{Règle~:} l'absence d'une trace exigée par le critère est un \NON{}
établi. \NEK{} est réservé au cas où l'auditeur ne peut pas savoir.
\end{boxalert}

\begin{boxscenario}{Dossier MINÉQUIR. Le contrôle 5 est levé}
\textbf{Constatation.} Le 2~avril~2026, le poste comptable remet les
visas électroniques des mandats 2025-08817 et 2025-09342. Les deux
portent la vérification des pièces jointes~; les pièces étaient
complètes et lisibles.

\medskip
\textbf{Analyse} (confiance élevée). Le contrôle des pièces a fonctionné
tel qu'il est conçu~: il vérifie la présence et la lisibilité des pièces,
pas leur unicité. Le contrôle~5 passe de \NEK{} à \OUI{}. Le double
paiement relève d'un autre contrôle, le blocage des doublons (contrôle~4),
coté \NON{}.
\end{boxscenario}

\separateur

% ============================================================================
\section{Classer les déficiences}
\label{sec:c03-deficiences}
% ============================================================================

\subsection{Deux niveaux}

Le COSO~2013 distingue deux niveaux de déficience~\cite{coso2013icfr}.

\begin{boxdef}
\termecle{Déficience du contrôle interne}\index{déficience}~: insuffisance
d'une ou plusieurs composantes ou d'un ou plusieurs principes, qui réduit
la probabilité que l'entité atteigne ses objectifs.

\medskip
\termecle{Déficience majeure}\index{déficience majeure}~: déficience, ou
combinaison de déficiences, qui réduit \emph{fortement} cette probabilité.
\end{boxdef}
{\small\itshape Définitions reformulées d'après le COSO~2013, §~138.}

\medskip

Deux règles du COSO (§§~139 et~140) commandent la conclusion. Lorsqu'une
déficience majeure existe, l'entité ne peut pas conclure que son système
de contrôle interne est efficace. Et une déficience majeure d'une
composante, ou d'un principe, ne peut pas être compensée par le bon
fonctionnement d'une autre composante, ou d'autres principes.

\subsection{Une grille de classement}

Pour décider si une déficience est majeure, le cours retient quatre
questions. Elles ne remplacent pas le jugement~; elles l'obligent à
s'expliquer.

\begin{tableaucours}{cL{6.4cm}L{5.4cm}}{%
    \tableauHeader{N°} &
    \tableauHeader{Question} &
    \tableauHeader{Penche vers «~majeure~» si…}
}
1 & Quelle perte, ou quelle irrégularité, la déficience rend-elle
possible~? & Elle porte sur des montants importants ou sur tout le
circuit \\
2 & Avec quelle probabilité~? & Rien n'empêche qu'elle se produise, ou
elle s'est déjà produite \\
3 & Un contrôle compensatoire la couvre-t-il, et est-il prouvé~? & Aucun,
ou il n'est pas prouvé \\
4 & Ouvre-t-elle une possibilité de fraude~? & Oui, en particulier par
contournement de la séparation des tâches \\
\end{tableaucours}
\vspace{-0.8em}
\begin{center}{\small\itshape Tableau~3.7. Grille de classement des
déficiences (convention du cours)}\end{center}

\subsection{Les déficiences du MINÉQUIR}

Les cinq \NON{} de la matrice ne font pas cinq déficiences~: plusieurs
contrôles manquants peuvent relever d'une même cause. La vérificatrice
les regroupe en trois déficiences.

\begin{tableaucours}{cL{4.6cm}L{1.8cm}L{1.8cm}L{3.0cm}}{%
    \tableauHeader{} &
    \tableauHeader{Déficience} &
    \tableauHeader{Contrôles} &
    \tableauHeader{Principes} &
    \tableauHeader{Classement}
}
A & Les accès ne sont pas maîtrisés~: comptes des partants actifs,
séparation LIQ / ORD non imposée & 1, 2 & 10, 11 & Majeure (questions 1
à 4) \\
B & Le fichier des fournisseurs n'est pas protégé et les doublons ne sont
pas bloqués & 3, 4 & 10 & Déficience \\
C & Aucune revue des opérations de fin de gestion & 7 & 16 & Déficience \\
\end{tableaucours}
\vspace{-0.8em}
\begin{center}{\small\itshape Tableau~3.8. Classement des déficiences du
circuit de la dépense}\end{center}

\begin{boxinfo}
\textbf{Pourquoi B n'est pas majeure à elle seule.} Si les accès étaient
maîtrisés, une modification de RIB et un mandat en double exigeraient
deux personnes différentes, et la collusion resterait nécessaire. C'est
la déficience~A qui rend B dangereuse. Le classement dit cette relation~:
la combinaison A et B est majeure, et c'est elle qui a permis le double
paiement.
\end{boxinfo}

\separateur

% ============================================================================
\section{Conclure sur l'efficacité du système}
\label{sec:c03-conclure}
% ============================================================================

\subsection{La lecture de conformité}

\begin{boxlecture}{Circuit de la dépense, après les tests}
\textbf{Constatation.} Sur huit contrôles attendus, trois existent et
fonctionnent (contrôles~5, 6 et~8), cinq manquent ou n'ont pas fonctionné.
Aucun contrôle ne reste à établir.

\medskip
\textbf{Analyse} (confiance élevée). Les contrôles qui fonctionnent sont
la vérification des pièces et deux contrôles de détection. La déficience
majeure~A porte sur les accès, dont dépendent tous les autres contrôles.

\medskip
\hfill\conformite{3}{8}{0}
\end{boxlecture}

\subsection{Une échelle simple}

Pour situer un domaine et suivre son évolution d'une mission à l'autre,
le cours propose une échelle à cinq niveaux. C'est une convention
pédagogique, pas une norme~: elle aide à dire où l'on en est, elle ne
remplace pas la conclusion.

\begin{tableaucours}{cL{2.6cm}L{9.2cm}}{%
    \tableauHeader{Niveau} &
    \tableauHeader{Nom} &
    \tableauHeader{Description}
}
0 & Inexistant & Le risque n'est pas identifié~; aucun contrôle \\
1 & Informel & Des contrôles existent, par habitude, sans procédure ni
trace \\
2 & Défini & Les contrôles sont écrits, mais leur exécution n'est pas
prouvée \\
3 & Appliqué & Les contrôles sont écrits, exécutés et tracés~; les tests
de fonctionnement sont concluants \\
4 & Piloté & Les contrôles sont surveillés, les déficiences suivies et
corrigées \\
\end{tableaucours}
\vspace{-0.8em}
\begin{center}{\small\itshape Tableau~3.9. Échelle de maturité du
contrôle interne (convention du cours)}\end{center}

Le circuit de la dépense du MINÉQUIR est au niveau~1~: la séparation des
tâches existe par habitude (cheminement), sans être imposée ni écrite.

\subsection{Formuler la conclusion}

La conclusion porte sur un objectif, une période et un périmètre, et dit
sur quoi elle repose. Celle de Mme~NGO~BIYONG s'écrit ainsi~:

\begin{quote}\itshape
Sur la base des travaux décrits au rapport (matrice des risques et des
contrôles, questionnaire, test de cheminement, tests de fonctionnement
sur l'exercice 2025), le système de contrôle interne du circuit de la
dépense de la DAF du MINÉQUIR n'est pas efficace~: il présente une
déficience majeure dans la maîtrise des accès à SIDÉPENSE, qui a permis
un double paiement de 187\,450\,000~FCFA.
\end{quote}

\begin{boxprocedure}
\textbf{Le dossier d'évaluation d'un circuit} contient, dans cet ordre~:
\begin{enumerate}
    \item le narratif et le diagramme, validés par l'entité~;
    \item la matrice des risques et des contrôles, datée~;
    \item le QCI rempli, avec le nom et la fonction de chaque répondant~;
    \item la fiche du test de cheminement et ses traces~;
    \item une feuille de test par contrôle~: conception, échantillon,
          résultat, cote, pièces~;
    \item le tableau de classement des déficiences, avec la réponse aux
          quatre questions pour chaque déficience majeure~;
    \item la conclusion, et la référence des pièces sur lesquelles elle
          repose.
\end{enumerate}
\end{boxprocedure}

\begin{boxmetier}{Vérificateur}
L'article~28 du décret n°~2013/287 vous charge d'évaluer les structures
de contrôle interne. Joignez à votre rapport la matrice des risques et des
contrôles~: elle montre que votre conclusion repose sur une liste établie
à l'avance, et non sur ce que l'entité a bien voulu montrer.
\end{boxmetier}

\begin{boxmetier}{Inspecteur}
Une inspection déclenchée par une alerte peut sauter l'étape du QCI, mais
pas celle du cheminement~: suivez au moins une opération réelle dans
l'application avant de conclure sur un circuit.
\end{boxmetier}

\begin{boxmetier}{Auditeur interne}
Votre évaluation annuelle du contrôle interne sera relue par le CONSUPE.
Tenez une matrice des risques et des contrôles par circuit, et mettez-la
à jour à chaque changement d'application~: c'est votre meilleure défense
devant un vérificateur.
\end{boxmetier}

\separateur

\begin{boxresume}
\begin{itemize}
    \item \textbf{Cinq étapes.} Décrire, questionner, cheminer, évaluer la
          conception, tester le fonctionnement~; puis classer et conclure.
    \item \textbf{La matrice part des risques}, pas des contrôles déclarés.
    \item \textbf{Le QCI recueille des déclarations}~; le cheminement et les
          tests produisent des preuves.
    \item \textbf{Conception d'abord.} Un contrôle mal conçu est une
          déficience, sans test de fonctionnement.
    \item \textbf{Trois cotes.} \OUI{}, \NON{}, \NEK{}~; l'absence d'une
          trace exigée est un \NON{}.
    \item \textbf{Deux niveaux de déficience.} Une déficience majeure
          interdit de conclure à l'efficacité, et ne se compense pas.
    \item \textbf{Une conclusion située.} Un objectif, une période, un
          périmètre, et les travaux sur lesquels elle repose.
\end{itemize}
\end{boxresume}

\begin{boxexo}[title={\ding{45}\quad Exercice~3.1. Formuler des questions \niveaubase}]
Réécrivez les questions suivantes du QCI pour qu'elles appellent une
réponse vérifiable~: (a) «~Les sauvegardes sont-elles faites~?~»~; (b)
«~Les mots de passe sont-ils sûrs~?~»~; (c) «~Le fichier de paie est-il
contrôlé~?~».
\end{boxexo}

\begin{boxexo}[title={\ding{45}\quad Exercice~3.2. Dimensionner et coter \niveaubase}]
Le chef de la liquidation signe chaque soir la liste des mandats du jour.
\begin{enumerate}[(a)]
    \item Combien de listes testez-vous pour l'exercice 2025~?
    \item Sur l'échantillon, deux listes ne sont pas signées. Quelle
          cote~?
    \item Les listes de novembre n'ont pas été retrouvées, et le service
          affirme qu'elles ont été signées. Quelle cote pour novembre, et
          que faites-vous~?
\end{enumerate}
\end{boxexo}

\begin{boxexo}[title={\ding{45}\quad Exercice~3.3. Une matrice pour la paie \niveaumoyen}]
Construisez la matrice des risques et des contrôles du circuit de la paie
de la CAPNO (entrée d'un agent, modification d'un élément de solde,
départ, calcul, virement), avec au moins six risques, le contrôle attendu
pour chacun, son type et le principe du COSO correspondant.
\end{boxexo}

\begin{boxexo}[title={\ding{45}\quad Exercice~3.4. Classer et conclure \niveauexpert}]
À la CAPNO, l'auditeur interne relève~: (1) 37 agents payés ne figurent
pas au fichier du personnel~; (2) le profil qui crée un agent peut aussi
valider la paie~; (3) aucun rapprochement entre la paie et le fichier du
personnel n'est fait~; (4) les sauvegardes ne sont pas testées.
\begin{enumerate}[(a)]
    \item Regroupez ces constats en déficiences, et classez chacune à
          l'aide du Tableau~3.7.
    \item Rédigez la conclusion sur l'efficacité du contrôle interne de la
          paie.
    \item Dites ce qui changerait votre classement.
\end{enumerate}
\end{boxexo}

\begin{boxinfo}
\textbf{Transition.} La Partie~I a donné l'objet et la méthode~:
évaluer le contrôle interne d'un système d'information. Reste à savoir au
nom de quoi l'on évalue. Le Chapitre~\ref{chap:institutions} présente les
institutions qui contrôlent les systèmes d'information publics au
Cameroun, et leurs compétences.
\end{boxinfo}
